%%======================================================================
\section{Discussion and Future Work}
\label{sec:discussion}
%%======================================================================

The evidence surveyed in Sections~\ref{sec:applications} and~\ref{sec:challenges}
shows a field that can demonstrate clinical relevance in laboratory settings but
has not cleared any of the three thresholds required for operational deployment:
accepted safety standards, clinical evidence in human subjects, and regulatory
precedent.

\subsection{Technology Readiness Assessment}
\label{sec:trl}

The TRL assignments in Table~\ref{tab:evidence} provide the starting point. The
aggregate finding (TRL~3--4 across all domains for bipedal humanoid
platforms) is not a uniform statement; the TRL ceiling for each domain is
governed by a specific domain-particular barrier.

\textbf{Clinical procedures (current TRL~3--4).} Atar et
al.~\cite{atar2025humanoids}, Liang et al.~\cite{liang2025lapsurgie}, and Cho et
al.~\cite{cho2026humanoid} constitute the most technically substantive evidence
in the review. The cadaveric study of Cho et al.\ is the current frontier of the
evidence sequence: the step immediately preceding an IRB-approved human-subject
feasibility trial. Advancing to TRL~5 requires two things simultaneously: a
structured IRB-approved feasibility study in a specific narrow indication
(telemedicine-assisted physical examination is the most defensible starting point,
because fully teleoperated operation keeps a human in the decision loop at all
times), and independent replication on a platform other than the Unitree G1.

\textbf{Elderly and nursing care (TRL~4--5 for social robots; TRL~3--4 for
bipedal humanoids).} The gap between platform classes is analytically significant.
The SPRING project achieved TRL~4--5 through real-world deployment with 60+~older
adult users over four years~\cite{alamedapineda2024spring}. RHP
Friends, the only purpose-built bipedal nursing humanoid, demonstrated patient
transfer tasks at IREX~2023 but has not been evaluated in an actual care
facility~\cite{benallegue2025rhp}. The path to TRL~5 for bipedal humanoids is
not a question of platform capability but of deployment context: a structured
pilot in an actual care facility, using SPRING-derived methodology on a bipedal
platform, is the specific experiment that would advance this domain.

\textbf{Rehabilitation and mental health (TRL~3--4 and TRL~4--5, respectively).}
These two domains share a structural problem: the strongest available evidence
was generated on the wrong platforms. In rehabilitation, Nguyen et
al.~\cite{nguyen2024eeg} establish neuroscientific rationale, Sobrepera et
al.~\cite{sobrepera2022feasibility} establish clinical acceptability in a small
case series, and Lu et al.~\cite{lu2025gentlehumanoid} establish the
contact-safety engineering foundation, but none of the three used a full-size
modern bipedal humanoid with a clinical patient population. In mental health,
Robinson et al.~\cite{robinson2023brief} ($n\,=\,230$ RCT) and Sayis and
Gunes~\cite{sayis2024technology} ($n\,=\,42$) achieve methodologically strong
results with consistent embodiment advantage findings, but neither specifies the
humanoid model, and all participants are non-clinical wellness populations.

The advance required in both domains is the same experiment in different forms:
an IRB-approved human-subject study using a specified full-size bipedal platform
with a defined patient population. For rehabilitation, the target is a stroke
trial with standardized outcome measures (Fugl-Meyer Assessment,
10-meter walk test). For mental health, the target is a feasibility trial with a
clinically diagnosed population assessing validated scales. Both domains are one
targeted study away from TRL~5. Both are currently blocked by the absence of
that study, not by lack of underlying evidence.

\textbf{Hospital logistics (TRL~2).} No peer-reviewed evidence exists for
humanoid-specific hospital logistics tasks. Wheeled autonomous mobile robots
already address this domain with demonstrated operational reliability. Hospital
logistics is not a priority for bipedal humanoid development at this stage.

\textbf{Cross-domain finding.} No domain has reached TRL~7 (system prototype
demonstrated in operational environment). TRL~5 (the threshold of first
human-subject validation in a controlled clinical setting) is the realistic
five-year target for the three leading domains.

Figure~\ref{fig:trl} visualizes the current TRL state and TRL~5 target across
all five clinical domains.

\begin{figure}[t]
  \centering
  \Description{Horizontal bar chart showing the Technology Readiness Level (TRL)
  assessment for five clinical domains: clinical procedures (TRL 3--4), elderly
  and nursing care (TRL 3--5), rehabilitation (TRL 3--4), mental health and HRI
  (TRL 3--5), and hospital logistics (TRL 2). A dashed vertical line marks TRL 5
  as the five-year target for leading domains.}
  \includegraphics[width=0.75\columnwidth]{figures/trl_readiness.pdf}
  \caption{Technology Readiness Level (TRL) assessment by clinical domain
  (March 2026). Current TRL range (filled bars) and TRL~5 target (dashed line)
  per ISO 16290:2013. No domain has reached TRL~7 (operational environment
  prototype); TRL~5 (relevant-environment validation) is the realistic five-year
  target for clinical procedures, elderly care (bipedal), and rehabilitation.}
  \label{fig:trl}
\end{figure}
Figure~\ref{fig:capability_gap} quantifies the gap between best current commercial
performance and the scores required to reach TRL~5 for clinical deployment.

\begin{figure}[htbp]
  \centering
  \Description{Horizontal grouped bar chart with six rows (capability dimensions)
  and two bars per row. Blue bars show current scores (0-10 scale) for each
  dimension; orange bars show required scores for TRL 5 clinical deployment. The
  six dimensions are: Manipulation Precision (current 4, required 8), Bipedal
  Gait Stability (current 5, required 7), LLM/VLA Task Generalization (current 5,
  required 7), Contact Safety (current 3, required 8), Regulatory Compliance
  Readiness (current 2, required 8), and Human Trust Score (current 5, required 7).
  The largest gaps are in Regulatory Compliance Readiness and Contact Safety.}
  \includegraphics[width=\columnwidth]{figures/capability_gap.pdf}
  \caption{Humanoid capability gap: current best-commercial scores (blue, 2026)
  versus required scores for TRL~5 clinical deployment (orange). Scale 0--10
  represents engineering and evidence maturity. Regulatory Compliance Readiness
  and Contact Safety show the largest absolute gaps; Manipulation Precision is
  the dominant procedural engineering gap.}
  \label{fig:capability_gap}
\end{figure}

\subsection{Near-Term Deployable Roles}
\label{sec:roles}

Three roles have the strongest case for TRL~5 achievement within five years.

\textbf{Role~1: Telemedicine and telepresence clinical assistant.} The teleoperated
mode is the key risk-mitigation design choice: by keeping a human operator in the
decision loop at all times, it avoids the regulatory classification and safety
validation burdens associated with autonomous patient contact. The technical
infrastructure exists: Atar et al.~\cite{atar2025humanoids} demonstrate
teleoperated clinical task execution across seven procedures, and Liang et
al.~\cite{liang2025lapsurgie} extend this to laparoscopic surgery. The primary
application is specialist access extension: enabling a remote specialist to
conduct a physical examination or assist in a procedure in a rural or low-resource
setting. Realizing this role at TRL~5 requires a defined physical examination
protocol, remote operator training standards, and IRB oversight.

\textbf{Role~2: Rehabilitation movement demonstrator.} This role requires no
direct patient contact: the humanoid demonstrates movements that the patient
observes or follows. The bipedal gait safety gap identified in
Section~\ref{sec:safety}, the most universally blocking technical
barrier, becomes irrelevant to this specific role, because the safety concern is
dynamic contact force during locomotion near patients, not movement observation at
distance. The neuroscientific rationale is grounded in Nguyen et al.'s mirror
neuron evidence~\cite{nguyen2024eeg}; the clinical acceptability is established by
Sobrepera et al.'s case series~\cite{sobrepera2022feasibility}. Advancing requires a registered
feasibility trial adapting a NAO or Pepper protocol to a full-size bipedal
humanoid demonstrator in stroke rehabilitation.

\textbf{Role~3: Structured eldercare companion and physical assistant.} RHP
Friends~\cite{benallegue2025rhp} is the only platform designed specifically for
this role. The hybrid autonomous/teleoperated architecture (in which routine
tasks run autonomously and non-routine tasks are handed to a remote
operator) is the correct design for this deployment scenario. The SPRING
project~\cite{alamedapineda2024spring} provides the deployment methodology and
user-acceptance evaluation framework. Realizing this role at TRL~5 requires an
actual care facility pilot.

\subsection{Research Priorities}
\label{sec:research}

Four priorities, ordered by their blocking impact on the entire field:

\textbf{Priority~1: Clinical safety validation for bipedal gait in patient-contact
environments.} The four gait-induced contact risk scenarios documented in
Section~\ref{sec:safety} are not bounded by any existing standard. No IRB will
approve human-subject studies involving physical patient contact by a full-size
bipedal humanoid without instrumented biomechanical data on these force profiles
and an accepted safety standard or architectural safeguard. Among the priorities
identified here, this is the one most likely to affect multiple clinical domains,
because it would provide a basis for patient-contact studies across the leading
applications.

\textbf{Priority~2: Human-subject feasibility studies in leading applications.}
The cadaveric endoscopy study of Cho et al.~\cite{cho2026humanoid} is the current
clinical frontier. The specific next step is an IRB-approved, registered
feasibility study in a narrow indication with an explicit safety monitor protocol.
UC San Diego (Yip Laboratory) and Johns Hopkins are especially well positioned
for such work, given their existing contributions to the clinical procedures
literature.

\textbf{Priority~3: Open clinical demonstration dataset for VLA training.} No
public dataset of humanoid robot demonstrations in clinical or
clinical-simulated environments exists. Creating an open clinical demonstration
dataset, even using simulators and clinical task protocols with mannequins, would
unlock foundation model fine-tuning for clinical applications and enable the
broader research community to advance clinical policy learning.

\textbf{Priority~4: Platform diversification for evidence generalizability.} The
concentration of clinical procedure evidence on the Unitree G1 (Atar et al., Liang et al., Cho et al.) is a
scientific limitation. Replication of clinical task evaluations on RHP Friends,
TALOS, or other platforms would establish whether current findings are
platform-specific or generalizable to the bipedal humanoid class.

\subsection{Humanoid Clinical Deployment Readiness Framework}
\label{sec:hcdr}

The TRL scale, originally defined by NASA for space systems~\cite{mankins1995trl}
and later formalized as ISO~16290:2013~\cite{iso16290} for broader
application~\cite{heder2017trl}, captures technology maturity along a single
axis. Clinical deployment decisions require reasoning along at least three
independent axes simultaneously: \emph{autonomy level}, \emph{patient-contact intensity}, and
\emph{regulatory complexity}. We introduce the Humanoid Clinical Deployment
Readiness (HCDR) framework to make this three-dimensional space explicit. The
framework is not a scoring system; it is a classification tool for locating
specific role-platform combinations in the deployment readiness space and
identifying which axis is the binding constraint.

\textbf{Axis~1: Autonomy level.} Roles range from fully teleoperated
(human operator controls every action in real time) through supervised autonomy
(the robot executes programmed sub-tasks under human oversight and with human
abort authority) to conditional autonomy (robot operates autonomously under
defined operational conditions with human escalation) and full autonomy
(unsupervised operation). Current evidence supports only fully teleoperated roles.
No qualifying study demonstrates supervised autonomy in a patient-contact task.

\textbf{Axis~2: Patient-contact intensity.} Contact ranges from no direct contact
(observation, companionship, logistics) through indirect contact (physical
interaction via tools or instruments, patient not in direct force chain) to direct
physical contact (hands on patient, transfer assistance, examination). The binding
regulatory and safety constraints scale sharply with contact intensity: indirect
tool contact requires force control validation; direct physical contact requires
validated safety standards that do not yet exist for bipedal platforms.

\textbf{Axis~3: Regulatory complexity.} Complexity is determined by device
classification (Class~I--III under FDA; Class~I--III under EU MDR), presence of
AI/ML as a decision-making component, and whether the application involves
therapeutic intent (triggering clinical investigation requirements) or supportive
logistics (lower classification pathway). Full autonomy with direct patient contact
and therapeutic intent is the highest-complexity combination; teleoperation for
physical examination assistance with physician oversight is among the lowest.

Table~\ref{tab:hcdr} maps the three near-term deployable roles against the HCDR
framework and identifies the binding constraint on each path.

\begin{table}[t]
  \caption{Humanoid Clinical Deployment Readiness (HCDR) framework applied to
  near-term deployable roles (March 2026). Binding constraint is the axis whose
  gap most delays deployment. TRL = Technology Readiness Level (ISO 16290:2013).}
  \label{tab:hcdr}
  \footnotesize
  \begin{tabular}{@{}p{2.2cm}p{1.5cm}p{1.5cm}p{1.7cm}p{1.7cm}@{}}
    \toprule
    \textbf{Role} & \textbf{Autonomy} & \textbf{Contact} & \textbf{Reg.\ Complexity} & \textbf{Binding Constraint} \\
    \midrule
    Telemedicine assistant & Teleoperated & Indirect (tools) & Class~II / De~Novo & Safety standard: bipedal gait \\[4pt]
    Rehabilitation demonstrator & Supervised & None & Class~I--II & Evidence: patient trial \\[4pt]
    Eldercare companion / assist & Supervised-to-conditional & Direct & Class~IIb--III & Both: gait safety + contact standard \\[4pt]
    Laparoscopic surgical assistant & Teleoperated & Indirect (instruments) & Class~III & RCM precision + regulatory precedent \\
    \bottomrule
  \end{tabular}
\end{table}

The framework reveals that the binding constraint is not uniform across roles.
For the telemedicine assistant, the primary block is a missing safety standard for
bipedal gait in patient-proximate environments, an engineering and standards-body
problem, not a clinical evidence problem. For the rehabilitation demonstrator, the
block is clinical evidence: the safety case is straightforward (no direct contact),
but no powered clinical trial exists. For eldercare physical assistance, both gait
safety and direct-contact force standards are missing simultaneously. The HCDR
matrix operationalizes this distinction: investment in gait safety standards
unblocks the telemedicine and eldercare roles; investment in a rehabilitation
trial can proceed independently without waiting for standards resolution.

\subsection{Practitioner Guidance}
\label{sec:practitioners}

The evidence synthesized in this review has different implications for different
practitioner audiences.

\textbf{For clinical engineers and hospital technology leadership.} The commercial
humanoid platforms entering the market in 2024--2026 are designed for industrial
labor, not healthcare. Procurement decisions made in 2026--2027 should distinguish
between research-access platforms (Unitree G1: open SDK, lowest cost, only
platform with published clinical task evidence, but physically limited) and future
clinical platforms (RHP Friends: purpose-built but academic/prototype-only; Sanctuary
Phoenix Gen~8: highest haptic sensing specification but closed ecosystem). No
commercial platform is ready for unsupervised clinical use in any application.
Any hospital-based evaluation program should structure the initial phase as an IRB-
approved research protocol, not an operational pilot, to ensure safety monitoring,
data collection, and liability clarity. Partnering with an academic medical center
(not a robotics company) to co-design the evaluation protocol is the appropriate
model.

\textbf{For robotics engineers entering the healthcare domain.} The most impactful
contribution from the engineering community is not yet another manipulation
algorithm or VLA architecture, which is not where the field is currently
blocked. Four specific
infrastructure gaps documented in Section~\ref{sec:accessibility} represent
engineering problems with no current open-source solutions: (1) a runtime safety
layer enforcing ISO/TS~15066 contact force limits for bipedal dynamic scenarios;
(2) clinical workflow integration middleware connecting humanoid action APIs to
HL7 FHIR and hospital EHR systems; (3) a public clinical demonstration dataset
for VLA fine-tuning on medical procedures; and (4) validated medical simulation
environments with patient anatomy proxies and clinical instrument physics. Any
of these would materially advance the field more than platform capability
improvements at current TRL levels.

\textbf{For clinical researchers and trial designers.} Three study designs are
currently feasible and would individually advance the field: (1) A registered,
IRB-approved feasibility study of teleoperated humanoid physical examination
assistance in a telemedicine-enabled rural clinic or outpatient setting, the
lowest-complexity entry into patient-contact evidence. (2) A pilot RCT of
humanoid-delivered action observation therapy for stroke rehabilitation, adapting
the Sobrepera et al.\ protocol to a full-size bipedal platform and powered for
a functional outcome measure (Fugl-Meyer Assessment). (3) A prospective
multi-session care facility pilot adapting SPRING methodology to RHP Friends or
an equivalent purpose-built platform, measuring care quality and staff burden
outcomes. All three designs require a pre-specified safety monitoring protocol
and an a priori definition of the study-stopping criteria for unexpected adverse
events.

\subsection{Regulatory Roadmap}
\label{sec:roadmap}

Figure~\ref{fig:regulatory} illustrates the parallel regulatory tracks and their
key decision points for both jurisdictions.

\begin{figure}[htbp]
  \centering
  \Description{Two-column flowchart comparing the US FDA and EU regulatory
  pathways for a clinical humanoid robot. Left column (FDA): novel device with
  no predicate leads to De Novo pathway (21 CFR 513(f)(2)), Pre-Submission
  meeting, De Novo request submission with bench and clinical data, FDA review
  approximately 12 months, then Class II market authorization with post-market
  surveillance. Right column (EU): EU MDR Class IIb/III assessment with Notified
  Body, EU AI Act high-risk AI system conformity assessment (effective August
  2026), clinical investigation under MDR Article 62, CE marking plus AI Act
  declaration of conformity, then EU market authorization with post-market
  clinical follow-up. A red banner at the bottom highlights the critical gap in
  both tracks: no accepted safety standard exists for bipedal gait in
  patient-proximate environments, as ISO 13482 excludes medical devices and
  ISO/TS 15066 is designed for fixed-base collaborative arms.}
  \includegraphics[width=0.85\columnwidth]{figures/regulatory_pathways.pdf}
  \caption{Regulatory pathways for clinical humanoid robots: FDA De~Novo (US,
  left) and EU MDR plus EU AI Act (EU, right). Both tracks require the same
  foundational evidence (safety standards for bipedal gait in patient-proximate
  environments) that does not yet exist (bottom banner). FDA De~Novo is
  identified as the most plausible first-clearance pathway; EU MDR plus AI Act
  dual compliance represents the EU parallel track.}
  \label{fig:regulatory}
\end{figure}

The path to the first regulatory clearance requires three sequential steps.

\textbf{Step~1: Performance standards definition.} A joint effort among humanoid
manufacturers, academic medical centers (UCSD, JHU), and FDA's Center for Devices
and Radiological Health (CDRH) is required to define the Special Controls that
any De~Novo submission must meet: fall prevention performance requirements, patient
contact force limits for dynamic bipedal scenarios, teleoperation latency bounds,
and AI/ML software validation methodology consistent with FDA's Good Machine
Learning Practice principles and the December~2024 PCCP guidance.

\textbf{Step~2: IRB feasibility study generating clinical evidence.} A registered,
IRB-approved feasibility study in a specific narrow indication generates the
clinical evidence required for a De~Novo submission. This study must be designed
from the outset with regulatory submission in mind: pre-specified safety endpoints
and a defined safety monitor protocol.

\textbf{Step~3: De~Novo submission establishing a regulatory predicate.} The
De~Novo application proposes a new product code for bipedal humanoid clinical
assistants with Special Controls from Steps~1 and~2. First clearance establishes
a predicate device, enabling subsequent 510(k) applications.

For the EU, the parallel path requires: an ISO/TC~299 scope revision of
ISO~13482 to include clinical humanoids; EU MDR Class~IIa conformity assessment;
and EU AI Act compliance under high-risk AI obligations (effective August~2026 for
new AI systems; August~2027 for AI embedded in regulated medical products).
Regulatory precedents for managing layered framework compliance exist in the
surgical robotics domain.

\subsection{Ethical Considerations}
\label{sec:ethics}

\textbf{Dignity of care.} Physical care interactions require that the patient
feels cared for, not processed. The SPRING project finding that participants cited
eye-level conversational positioning as important for dignified
interaction~\cite{alamedapineda2024spring} provides direct empirical evidence that
form factor affects perceived dignity of care. The uncanny valley phenomenon, the
negative affective response elicited by human-like but not fully human
entities~\cite{mori2012uncanny}, is particularly relevant for humanoid design in
care settings, where distress rather than curiosity is the cost of a poor
perceptual response. Care robotics designers must incorporate patient dignity
explicitly as a design criterion alongside technical performance metrics.

\textbf{Liability and informed consent.} No current healthcare liability
framework assigns accountability for harm caused by an autonomous or
semi-autonomous robotic action in a care setting. The ethical dimensions of robot
care for elderly individuals, including dignity, autonomy, and emotional
manipulation risks, require explicit frameworks before
deployment~\cite{sharkey2012granny}. Before any human-subject deployment, consent
processes for humanoid-assisted care must be defined. These require deliberate
engagement between legal, ethical, and clinical stakeholders before the first
human-subject study is launched.

\textbf{Data privacy.} A humanoid operating in a patient room integrates cameras,
microphones, depth sensors, physiological monitors, and, in future deployments,
integration with EHR systems, into a single mobile platform in one of the most
privacy-sensitive environments imaginable. No data governance framework specific
to clinical humanoid sensor data currently exists.

\textbf{Workforce framing.} The evidence does not support a displacement
narrative. No clinical role exists that a humanoid currently performs reliably
enough to replace a trained nurse or clinician. The evidence-supported framing is
task offloading: routine, physically demanding, or cognitively repetitive tasks
that do not require clinical judgment could be offloaded, freeing clinicians for
judgment-requiring activities. The administrator survey of Imtiaz and
Khan~\cite{imtiaz2024perceptions} shows that the perceived risk of reduced human
contact is a primary institutional barrier; communicating the task-offloading
framing accurately to institutional decision-makers is as important as resolving
the technical barriers.

